\newgeometry{margin=.5cm}
\centering
\section*{Flowchart: SQL-Abfrageprozess}
\pagestyle{empty}
\begin{figure}[H]
  \centering
  \begin{tikzpicture}[
      node distance=4mm,
      box/.style={rectangle, rounded corners=3pt, draw=#1!80, fill=#1!10,
          very thick, text width=75mm, minimum height=13mm, align=left, inner sep=6pt},
      decision/.style={diamond, aspect=2.5, draw=#1!80, fill=#1!10, text width= 30mm, very thick, align=center, inner sep=1pt},
      num/.style={circle, draw=#1!80, fill=#1!80, text=white,
          font=\bfseries\large, minimum size=8mm},
      arr/.style={->, >=stealth, thick},
      skip/.style={->, >=stealth, thick, rounded corners=5pt},
      small/.style={font=\small}
    ]

    \node[box=gray, minimum height=10mm, align=center] (start)
    {\faIcon{play} \textbf{Start}\\SQL-Abfrage beginnt};

    \node[box=blue, below=of start] (select)
    {\faIcon{columns} \lstinline|SELECT|\\Wählt die Spalten / Ausdrücke aus,\\die zurückgegeben werden.};

    \node[box=blue, below=of select] (from)
    {\faIcon{database} \lstinline|FROM|\\Bestimmt die Tabelle(n).};

    \node[box=cyan, dotted, right=14mm of from, text width=55mm] (join)
    {\faIcon{link} \lstinline|JOIN| (optional)\\Verbindet Tabellen anhand\\von Bedingungen.};

    \node[box=ForestGreen, dotted, below=of from] (where)
    {\faIcon{filter} \lstinline|WHERE| (optional)\\Filtert die Zeilen nach Bedingungen\\(vor der Gruppierung!)};

    \node[decision=orange, below=of where] (dgroup)
    {\textbf{Gruppierung}};

    \node[box=orange, dotted, below=9mm of dgroup] (group)
    {\faIcon{layer-group} \lstinline|GROUP BY| (optional)\\Gruppiert die Zeilen nach einer\\oder mehreren Spalten.};

    \node[box=orange, dotted, below=of group] (having)
    {\faIcon{filter} \lstinline|HAVING| (optional)\\Filtert Gruppen nach Bedingungen\\ nach der Gruppierung.};

    \node[decision=purple, below=of having] (dorder)
    {\textbf{Sortierung}};

    \node[box=purple, dotted, below=9mm of dorder] (order)
    {\faIcon{sort-amount-down} \lstinline|ORDER BY| (optional)\\Sortiert das Ergebnis nach einer\\oder mehreren Spalten.};

    \node[decision=red, below=of order] (dlimit)
    {\textbf{Begrenzung}};

    \node[box=red, dotted, below=9mm of dlimit] (limit)
    {\faIcon{list-ol} \lstinline|LIMIT| (optional)\\Beschränkt die Anzahl der\\zurückgegebenen Zeilen.};

    \node[box=teal, below=of limit, align=center] (result)
    {\faIcon{check-circle} \textbf{Ergebnis}\\Resultat der SQL-Abfrage};

    % Nummern
    \node[num=blue, left=8mm of select] {1};
    \node[num=blue, left=8mm of from] {2};
    \node[num=cyan, above right=-4mm and -4mm of join] {2a};
    \node[num=ForestGreen, left=8mm of where] {3};
    \node[num=orange, left=8mm of group] {4};
    \node[num=orange, left=8mm of having] {5};
    \node[num=purple, left=8mm of order] {6};
    \node[num=red, left=8mm of limit] {7};

    % Hauptfluss
    \draw[arr] (start) -- (select);
    \draw[arr] (select) -- (from);
    \draw[arr] (from) -- (where);
    \draw[arr, dashed] (from.east) -- (join.west);
    \draw[arr, dashed] (join.south) |- (where.east);
    \draw[arr] (where) -- (dgroup);
    \draw[arr] (dgroup) -- node[right, text=orange!80!black] {\textbf{Ja}} (group);
    \draw[arr] (group) -- (having);
    \draw[arr] (having) -- (dorder);
    \draw[arr] (dorder) -- node[right, text=purple!80!black] {\textbf{Ja}} (order);
    \draw[arr] (order) -- (dlimit);
    \draw[arr] (dlimit) -- node[right, text=red!80!black] {\textbf{Ja}} (limit);
    \draw[arr] (limit) -- (result);

    % Überspringen
    \draw[skip] (dgroup.east) -- ++(25mm,0)
    node[right, text=green!60!black] {\textbf{Nein}}
    |- (dorder.east);

    \draw[skip] (dorder.east) -- ++(25mm,0)
    node[right, text=green!60!black] {\textbf{Nein}}
    |- (dlimit.east);

    \draw[skip] (dlimit.east) -- ++(25mm,0)
    node[right, text=green!60!black] {\textbf{Nein}}
    |- (result.east);

  \end{tikzpicture}
\end{figure}

\vspace{1cm}

\section*{Flowchart: SQL-Befehle für Tabellenmanagement}
\begin{figure}[H]
  \centering
  \begin{tikzpicture}[
      node distance=4mm,
      box/.style={rectangle, rounded corners=3pt, draw=#1!80, fill=#1!10,
          very thick, text width=75mm, minimum height=13mm, align=left, inner sep=6pt},
      num/.style={circle, draw=#1!80, fill=#1!80, text=white,
          font=\bfseries\large, minimum size=8mm},
      arr/.style={->, >=stealth, thick}
    ]

    \node[box=gray, minimum height=10mm, align=center] (start)
    {\faIcon{play} \textbf{Start}\\Tabelle existiert noch nicht};

    \node[box=blue, below=of start] (create)
    {\faIcon{table} \lstinline|CREATE TABLE|\\Erstellt eine neue Tabelle mit\\definierten Spalten und Typen.};

    \node[box=ForestGreen, below=of create] (insert)
    {\faIcon{plus-circle} \lstinline|INSERT INTO|\\Fügt neue Zeilen (Datensätze)\\in die Tabelle ein.};

    \node[box=orange, below=of insert] (update)
    {\faIcon{pen} \lstinline|UPDATE|\\Ändert bestehende Daten\\in bestimmten Zeilen.};

    \node[box=cyan, below=of update] (alter)
    {\faIcon{edit} \lstinline|ALTER TABLE|\\Ändert die Struktur der Tabelle\\(z.B. neue Spalten).};

    \node[box=purple, below=of alter] (delete)
    {\faIcon{minus-circle} \lstinline|DELETE FROM|\\Löscht bestimmte Zeilen\\aus der Tabelle.};

    \node[box=red, below=of delete] (drop)
    {\faIcon{trash-alt} \lstinline|DROP TABLE|\\Löscht die gesamte Tabelle\\inklusive aller Daten.};

    \node[box=teal, below=of drop, align=center] (end)
    {\faIcon{check-circle} \textbf{Ende}\\Tabelle existiert nicht mehr};

    % Nummern
    \node[num=blue, left=8mm of create] {1};
    \node[num=ForestGreen, left=8mm of insert] {2};
    \node[num=orange, left=8mm of update] {3};
    \node[num=cyan, left=8mm of alter] {4};
    \node[num=purple, left=8mm of delete] {5};
    \node[num=red, left=8mm of drop] {6};

    % Hauptfluss
    \draw[arr] (start) -- (create);
    \draw[arr] (create) -- (insert);
    \draw[arr] (insert) -- (update);
    \draw[arr] (update) -- (alter);
    \draw[arr] (alter) -- (delete);
    \draw[arr] (delete) -- (drop);
    \draw[arr] (drop) -- (end);

  \end{tikzpicture}
\end{figure}

\restoregeometry